\documentclass[11pt,a4paper]{article}

\usepackage[T1]{fontenc}
\usepackage[utf8]{inputenc}
\usepackage[polish]{babel}
\usepackage{geometry}
\usepackage{hyperref}
\usepackage{enumitem}
\usepackage{listings}
\usepackage{xcolor}
\usepackage{titlesec}

\geometry{margin=2.5cm}

\hypersetup{
    colorlinks=true,
    linkcolor=blue,
    urlcolor=blue
}

\lstset{
    basicstyle=\ttfamily\small,
    breaklines=true,
    frame=single,
    columns=fullflexible
}

\titleformat{\section}{\large\bfseries}{\thesection.}{0.5em}{}
\titleformat{\subsection}{\normalsize\bfseries}{\thesubsection.}{0.5em}{}

\title{Instrukcja uruchamiania i sprawdzania działania backendu na Windows}
\author{}
\date{}

\begin{document}

\maketitle

\section{Cel dokumentu}

Ten dokument opisuje, jak uruchomić backend lokalnie na systemie Windows oraz jak krok po kroku sprawdzić, czy działa poprawnie. Instrukcja jest przeznaczona dla programistów, którzy dołączają do projektu i chcą szybko uruchomić środowisko developerskie na własnej maszynie.

\section{Wymagania wstępne}

Do uruchomienia projektu potrzebne są:

\begin{itemize}[leftmargin=1.2cm]
    \item Windows 10 lub nowszy,
    \item Python 3.9 lub nowszy,
    \item pip,
    \item PowerShell albo Command Prompt,
    \item dostęp do katalogu \texttt{backend/}.
\end{itemize}

Projekt działa lokalnie na bazie \textbf{SQLite}, więc nie trzeba osobno uruchamiać PostgreSQL.

\section{Struktura robocza}

Zakładamy, że projekt ma strukturę:

\begin{lstlisting}
project-root/
|-- backend/
|   |-- app/
|   |-- tests/
|   |-- requirements.txt
|   `-- run.py
|-- frontend/
`-- mobile/
\end{lstlisting}

Cała instrukcja dotyczy wyłącznie katalogu \texttt{backend/}.

\section{Uruchomienie projektu krok po kroku na Windows}

\subsection{1. Przejdź do katalogu backend}

W PowerShell:

\begin{lstlisting}
cd backend
\end{lstlisting}

\subsection{2. Utwórz środowisko wirtualne}

\begin{lstlisting}
python -m venv .venv
\end{lstlisting}

Jeżeli komenda \texttt{python} nie działa, spróbuj:

\begin{lstlisting}
py -m venv .venv
\end{lstlisting}

\subsection{3. Aktywuj środowisko}

W PowerShell:

\begin{lstlisting}
.venv\Scripts\Activate.ps1
\end{lstlisting}

W Command Prompt:

\begin{lstlisting}
.venv\Scripts\activate.bat
\end{lstlisting}

Po aktywacji terminal powinien pokazać prefiks \texttt{(.venv)}.

\subsection{4. Zainstaluj zależności}

\begin{lstlisting}
pip install -r requirements.txt
\end{lstlisting}

\subsection{5. Utwórz plik .env}

W katalogu \texttt{backend/} utwórz plik \texttt{.env} i wpisz:

\begin{lstlisting}
FLASK_ENV=development
DATABASE_URL=sqlite:///app.db
SECRET_KEY=change-me
JWT_SECRET_KEY=change-me-too
CORS_ORIGINS=http://localhost:3000
APP_TIMEZONE=Europe/Warsaw
MAIL_SENDER=noreply@example.com
MAIL_BACKEND=console
\end{lstlisting}

\subsection{6. Usuń starą bazę, jeśli chcesz wystartować od zera}

W PowerShell:

\begin{lstlisting}
Remove-Item app.db -ErrorAction SilentlyContinue
Remove-Item instance\app.db -ErrorAction SilentlyContinue
\end{lstlisting}

W Command Prompt:

\begin{lstlisting}
del app.db
del instance\app.db
\end{lstlisting}

Ten krok nie jest obowiązkowy, ale warto go wykonać przy pierwszym uruchomieniu lub po większych zmianach w seedzie.

\subsection{7. Uruchom backend}

\begin{lstlisting}
python run.py
\end{lstlisting}

Jeżeli \texttt{python} nie działa, użyj:

\begin{lstlisting}
py run.py
\end{lstlisting}

\subsection{8. Oczekiwany efekt}

W terminalu powinien pojawić się komunikat podobny do:

\begin{lstlisting}
* Serving Flask app 'app'
* Debug mode: on
* Running on http://127.0.0.1:5000
\end{lstlisting}

Oznacza to, że backend działa lokalnie na porcie \texttt{5000}.

\section{Co robi uruchomienie przez run.py}

Aktualnie \texttt{run.py} działa w trybie developerskim i wykonuje następujące kroki:

\begin{itemize}[leftmargin=1.2cm]
    \item tworzy aplikację Flask,
    \item inicjalizuje bazę danych,
    \item wykonuje \texttt{db.create\_all()},
    \item uruchamia seed danych testowych,
    \item startuje lokalny serwer HTTP.
\end{itemize}

To oznacza, że po uruchomieniu projektu baza SQLite zostaje przygotowana automatycznie i nie trzeba ręcznie wykonywać migracji.

\section{Konta testowe dostępne po seedzie}

Po starcie backendu dostępne są przykładowe konta:

\begin{itemize}[leftmargin=1.2cm]
    \item admin: \texttt{admin@example.com} / \texttt{Admin123!}
    \item terapeuta: \texttt{anna@example.com} / \texttt{Password123!}
    \item terapeuta: \texttt{piotr@example.com} / \texttt{Password123!}
    \item pacjent: \texttt{jan@example.com} / \texttt{Password123!}
    \item pacjent: \texttt{ola@example.com} / \texttt{Password123!}
\end{itemize}

\section{Jak sprawdzić, czy backend działa}

\subsection{1. Healthcheck}

Najprostszy test:

\begin{lstlisting}
curl http://127.0.0.1:5000/api/v1/health
\end{lstlisting}

Oczekiwany wynik:

\begin{lstlisting}
{
  "data": {
    "status": "ok"
  }
}
\end{lstlisting}

Jeżeli \texttt{curl} nie działa zgodnie z oczekiwaniem w PowerShell, możesz po prostu otworzyć w przeglądarce:

\begin{lstlisting}
http://127.0.0.1:5000/api/v1/health
\end{lstlisting}

\subsection{2. Lista usług}

\begin{lstlisting}
curl http://127.0.0.1:5000/api/v1/services
\end{lstlisting}

Oczekiwany wynik: odpowiedź JSON zawierająca listę usług w polu \texttt{data}.

\subsection{3. Logowanie}

\begin{lstlisting}
curl -X POST http://127.0.0.1:5000/api/v1/auth/login ^
  -H "Content-Type: application/json" ^
  -d "{\"email\":\"jan@example.com\",\"password\":\"Password123!\"}"
\end{lstlisting}

W PowerShell można też użyć:

\begin{lstlisting}
Invoke-RestMethod -Method POST `
  -Uri "http://127.0.0.1:5000/api/v1/auth/login" `
  -ContentType "application/json" `
  -Body '{"email":"jan@example.com","password":"Password123!"}'
\end{lstlisting}

Oczekiwany wynik:
\begin{itemize}[leftmargin=1.2cm]
    \item kod 200,
    \item \texttt{accessToken},
    \item \texttt{refreshToken},
    \item dane użytkownika.
\end{itemize}

\subsection{4. Endpoint wymagający autoryzacji}

Po zalogowaniu należy skopiować \texttt{accessToken} i wykonać:

\begin{lstlisting}
curl http://127.0.0.1:5000/api/v1/users/me ^
  -H "Authorization: Bearer <accessToken>"
\end{lstlisting}

W PowerShell:

\begin{lstlisting}
Invoke-RestMethod -Method GET `
  -Uri "http://127.0.0.1:5000/api/v1/users/me" `
  -Headers @{ Authorization = "Bearer <accessToken>" }
\end{lstlisting}

Oczekiwany wynik: dane aktualnie zalogowanego użytkownika.

\section{Przykładowy test przepływu pacjenta}

Poniżej znajduje się najprostszy scenariusz manualny, który warto przejść po uruchomieniu projektu.

\subsection{Krok 1. Zaloguj się}

Wykonaj request do:

\begin{lstlisting}
POST /api/v1/auth/login
\end{lstlisting}

i zapisz \texttt{accessToken}.

\subsection{Krok 2. Pobierz listę usług}

\begin{lstlisting}
GET /api/v1/services
\end{lstlisting}

Z odpowiedzi zapisz \texttt{service\_id}.

\subsection{Krok 3. Pobierz terapeutów dla usługi}

\begin{lstlisting}
GET /api/v1/services/<service_id>/therapists
\end{lstlisting}

Z odpowiedzi zapisz \texttt{therapist\_id}.

\subsection{Krok 4. Pobierz dostępność}

\begin{lstlisting}
GET /api/v1/availability?serviceId=...&therapistId=...&from=...&to=...
\end{lstlisting}

Z odpowiedzi wybierz pierwszy poprawny \texttt{startAt}.

\subsection{Krok 5. Zarezerwuj wizytę}

\begin{lstlisting}
curl -X POST http://127.0.0.1:5000/api/v1/appointments ^
  -H "Content-Type: application/json" ^
  -H "Authorization: Bearer <accessToken>" ^
  -d "{\"serviceId\":\"<service_id>\",\"therapistId\":\"<therapist_id>\",\"startAt\":\"<slot_startAt>\"}"
\end{lstlisting}

Oczekiwany wynik:
\begin{itemize}[leftmargin=1.2cm]
    \item kod 201,
    \item status wizyty \texttt{BOOKED}.
\end{itemize}

\subsection{Krok 6. Pobierz swoje wizyty}

\begin{lstlisting}
curl http://127.0.0.1:5000/api/v1/appointments/me ^
  -H "Authorization: Bearer <accessToken>"
\end{lstlisting}

Oczekiwany wynik: nowo utworzona wizyta znajduje się w liście.

\subsection{Krok 7. Anuluj wizytę}

\begin{lstlisting}
curl -X POST http://127.0.0.1:5000/api/v1/appointments/<appointment_id>/cancel ^
  -H "Content-Type: application/json" ^
  -H "Authorization: Bearer <accessToken>" ^
  -d "{\"reason\":\"Zmiana planów\"}"
\end{lstlisting}

Oczekiwany wynik:
\begin{itemize}[leftmargin=1.2cm]
    \item kod 200,
    \item status \texttt{CANCELLED\_BY\_PATIENT}.
\end{itemize}

\section{Jak uruchomić testy automatyczne na Windows}

\subsection{1. Upewnij się, że jesteś w katalogu backend}

\begin{lstlisting}
cd backend
\end{lstlisting}

\subsection{2. Uruchom testy}

Najbezpieczniej:

\begin{lstlisting}
python -m pytest tests -q
\end{lstlisting}

Jeżeli \texttt{python} nie działa, użyj:

\begin{lstlisting}
py -m pytest tests -q
\end{lstlisting}

Jeżeli środowisko wymaga jawnego ustawienia ścieżki modułów, użyj w PowerShell:

\begin{lstlisting}
$env:PYTHONPATH="."
python -m pytest tests -q
\end{lstlisting}

albo w Command Prompt:

\begin{lstlisting}
set PYTHONPATH=.
python -m pytest tests -q
\end{lstlisting}

\subsection{3. Co testują testy}

Testy obejmują głównie:
\begin{itemize}[leftmargin=1.2cm]
    \item auth,
    \item katalog usług,
    \item dostępność,
    \item rezerwacje i anulowanie,
    \item review/admin/staff.
\end{itemize}

\subsection{4. Jak działa środowisko testowe}

W testach:
\begin{itemize}[leftmargin=1.2cm]
    \item aplikacja uruchamiana jest w konfiguracji \texttt{testing},
    \item baza jest tworzona od nowa,
    \item seed danych ładowany jest automatycznie.
\end{itemize}

\section{Najczęstsze problemy i szybkie rozwiązania}

\subsection{Problem: ModuleNotFoundError: No module named 'app'}

Uruchom testy tak:

\begin{lstlisting}
python -m pytest tests -q
\end{lstlisting}

albo ustaw \texttt{PYTHONPATH}.

PowerShell:

\begin{lstlisting}
$env:PYTHONPATH="."
python -m pytest tests -q
\end{lstlisting}

Command Prompt:

\begin{lstlisting}
set PYTHONPATH=.
python -m pytest tests -q
\end{lstlisting}

\subsection{Problem: backend działa, ale dane są nieaktualne}

Usuń lokalną bazę SQLite i uruchom aplikację ponownie.

PowerShell:

\begin{lstlisting}
Remove-Item app.db -ErrorAction SilentlyContinue
Remove-Item instance\app.db -ErrorAction SilentlyContinue
python run.py
\end{lstlisting}

Command Prompt:

\begin{lstlisting}
del app.db
del instance\app.db
python run.py
\end{lstlisting}

\subsection{Problem: test anulowania wizyty zwraca 403 zamiast 200}

Najczęściej test wybrał slot, którego nie można już anulować zgodnie z regułą 24 godzin. W takim przypadku helper testowy powinien wybierać termin oddalony o więcej niż 24 godziny od bieżącego czasu.

\subsection{Problem: różnice między datetime aware i naive}

SQLite potrafi zwracać daty bez informacji o strefie. Jeżeli pojawia się błąd porównywania dat, należy ujednolicić obsługę czasu w serwisie lub dopiąć \texttt{timezone.utc} przed porównaniem.

\subsection{Problem: PowerShell blokuje aktywację .ps1}

Jeżeli przy aktywacji środowiska pojawia się błąd polityki wykonywania skryptów, uruchom PowerShell jako użytkownik i wykonaj:

\begin{lstlisting}
Set-ExecutionPolicy -Scope CurrentUser RemoteSigned
\end{lstlisting}

Następnie ponownie aktywuj środowisko:

\begin{lstlisting}
.venv\Scripts\Activate.ps1
\end{lstlisting}

\section{Zalecany sposób pracy dla kolejnych programistów}

Przy rozwijaniu backendu najlepiej pracować w tej kolejności:

\begin{enumerate}[leftmargin=1.2cm]
    \item doprecyzować model danych,
    \item przygotować lub poprawić schemat wejścia/wyjścia,
    \item zaimplementować logikę biznesową w serwisie,
    \item wystawić endpoint w \texttt{app/api/},
    \item dodać testy,
    \item zaktualizować seed, jeśli nowa funkcjonalność tego wymaga.
\end{enumerate}

To pozwala utrzymać porządek i nie rozlewać logiki biznesowej po route'ach.

\section{Krótka lista startowa dla nowego programisty}

Jeśli wchodzisz do projektu po raz pierwszy, wykonaj:

\begin{enumerate}[leftmargin=1.2cm]
    \item przejdź do katalogu \texttt{backend/},
    \item utwórz \texttt{.venv},
    \item aktywuj środowisko,
    \item zainstaluj zależności,
    \item utwórz plik \texttt{.env},
    \item usuń starą bazę SQLite,
    \item uruchom \texttt{python run.py},
    \item sprawdź \texttt{/api/v1/health},
    \item zaloguj się kontem z seeda,
    \item przejdź podstawowy flow pacjenta,
    \item uruchom \texttt{python -m pytest tests -q}.
\end{enumerate}

Po tych krokach środowisko powinno działać poprawnie i być gotowe do dalszego rozwoju.

\end{document}